\documentclass[11pt]{article}
\usepackage[a4paper,margin=21mm]{geometry}
\usepackage{fontspec}
\setmainfont{Latin Modern Roman}
\newfontfamily\CJKfont{Noto Serif CJK SC}
\XeTeXlinebreaklocale "zh"
\XeTeXlinebreakskip=0pt plus 1pt
\usepackage{amsmath,amssymb,amsthm,amscd,graphicx,color}
\usepackage{xurl}
\usepackage[unicode,hidelinks]{hyperref}
\allowdisplaybreaks
\setlength\emergencystretch{3em}
\numberwithin{equation}{section}
\newtheorem{Theorem}{Theorem}[section]
\newtheorem{Corollary}[Theorem]{Corollary}
\newtheorem{Lemma}[Theorem]{Lemma}
\newtheorem{Proposition}[Theorem]{Proposition}
 { \theoremstyle{definition}
\newtheorem{Definition}[Theorem]{Definition}
\newtheorem{Example}[Theorem]{Example}
\newtheorem{Remark}[Theorem]{Remark} }

\def\angd{120}
\def\dirp{{\slash \mkern -9mu p}}
\def\cI{{\mathcal I}}

\let\veps\varepsilon
\DeclareMathOperator{\id}{Id}
\DeclareMathOperator{\Tr}{Tr}

\def\nfrac#1#2{{\textstyle\frac{#1}{#2}}}

\begin{document}
\begin{center}\Large General numerator formula in graph-parametric integrals\end{center}
\noindent\textbf{Stable ID:} 550\quad\textbf{Author:} Marc P. Bellon\par
\noindent\textbf{Work:} Numerators in Parametric Representations of Feynman Diagrams\par
\noindent\textbf{Source:} \url{https://sigma-journal.com/2025/007/}\par
\noindent\textbf{Version:} Official SIGMA21(2025) journal PDF and TeX; acquired2026-09-30; exact bytes pinned by SHA256\par
\noindent\textbf{Rights:} CC BY 4.0. Authors retain copyright. \url{https://creativecommons.org/licenses/by/4.0/}. Reuse and typesetting adaptation; original language and mathematical content preserved.\par
\noindent\textbf{Difficulty (prior selection preserved):} {\CJKfont H2: 须将任意标量积的 Wick 配对按循环重组，精确处理重数与维数因子，再通过 Dodgson 行列式和 Stirling 消去证明图多项式公式。单个 Gaussian 积分或现成积分代入不足以完成该统一组合恒等式，全部循环重组论证原文明确给出。}\par
\medskip\noindent\textbf{Editorial boundary note (not author text):} Complete original main Theorem 2.2, the full author Gaussian-pairing/permutation averaging proof and partition/Stirling-coefficient conversion. Retains exact graph-matrix/minor sign conventions, completed-diagram reduction and scale-invariant projective setting, and the author Dodgson-identity appendix proof. Wick/Gaussian integration and Gamma functional equation remain standard cited background. The unrelated fermion-loop computations, graphical applications and later examples are excluded.\par
\medskip\hrule\medskip\noindent\textbf{Author text begins below.}\par
\setcounter{section}{1}
% Original source lines 168--250
\section{Presentation of the main results} \label{presentation}


The notations for graph polynomials will be the ones in the works of
Brown and collabora\-tors~\cite{Br2009,BrYe09}. To each connected graph~$G$ a matrix $\tilde M_G$ is associated.
The matrix $\tilde M_G$ is a block
matrix with a diagonal block with the Schwinger parameters $x_i$ associated to the edges as diagonal values, two
blocks made of the incidence matrix and minus its transpose and a zero block to complete it. Schematically, it can be written
\begin{equation*}
	\tilde M_G = \begin{pmatrix}
		\operatorname{Diag}(\mathbf{x}) & -{}^t \operatorname{Inc} \\ \operatorname{Inc} & 0
	\end{pmatrix}.
\end{equation*}
This matrix is singular since the sum over the rows indexed by the vertices of the incidence matrix is zero, but the
determinant of any matrix $M_G$ obtained by removing the row and column corresponding to any vertex gives the first
Symanzik polynomial of the graph~$U_G$, independently of the chosen vertex. It is further independent of any of the choices made in the
presentation of the matrix, like the order of the edges and vertices or the orientation of the edges. This polynomial is of degree~$L$, with
$L$ the dimension of the first homology group of the graph, the loop number in the language of physicists.

The objects which will be used to express our results are the polynomials
$U^{I ,J}_K$ defined for~$I$,~$J$ and~$K$ subsets of the edges of the graph
such that $I$ and~$J$ have the same number of elements~$k$. The polynomial
$U^{I ,J}_K$ is the determinant of the matrix deduced from the graph matrix $M_G$ by
removing the rows associated to $I$, the columns associated to $J$ and setting
the variables associated to $K$ to zero and has degree $L-k$. From the symmetry properties of $M_G$,
one deduces that $U^{I ,J}_K = U^{J,I}_K$. These graph polynomials may change sign but are otherwise independent of
the choices made in the presentation of the graph. Since
we will need a definite sign for these expressions, it will be preferable to define
$U^{I ,J}$ as the determinant of the matrix where the rows indexed by $I$ are
replaced with rows with only one non zero value 1 in a position corresponding to a row in
$J$. Expansion along the columns indexed by~$J$ or the rows indexed by~$I$ gives back the previous definition up to
a sign, but now the result only depends on the orientation of the edges indexed by
$I$ and $J $ and the bijection between these subsets
defined by the positions of the 1 in the modified matrix. Schnetz in~\cite{SC2021} proposed to put
an explicit sign in front of the determinant of the reduced matrix. My implicit definition turns out to be fully equivalent to his
explicit one. Both approaches can be useful.

These polynomials, in the special case where $I$ and $J$ have only one element,
give the adjoint elements of the matrix $M_G$, so that the correlation between momenta in rows
$i$ and $j$ is given by $U^{\{i\},\{j\}}/U=U^{i,j}/U$ (we do not keep the braces in the case of
singletons to lighten the notations). The Dodgson identities relate determinants of matrices whose elements
are $U^{i,j}$ for~${i \in I}$ and $j \in J$ to the polynomials $U^{I ,J}$. In Appendix \ref{polyn}, I state these
identities and their proofs, paying due attention to the question of signs.

Our results is most easily described for a vacuum diagram. The case of a propagator diagram can be reduced to this one, as will be seen in next
section.
The graph is replaced by the related completed graph,
obtained by linking the two exterior legs in a simple edge. In this case, the exponent associated to the added edge
must be chosen to make the whole vacuum diagram scale invariant.

We consider a numerator expressed as the product for $s=(s_0,s_1)$ in a set~$S$ of $n$ scalar
products $p_{s_0} \cdot p_{s_1}$. A priori, there are no restrictions to the number of times the momentum
associated to a given edge can appear in this product of scalar products. It should be even possible to have higher
powers of some of the scalar products, but it seems to enter the domain of useless generality.
\begin{Definition}
For a non-empty subset~$S_0$ of~$S$, define $U^{S_0}$ as the mean of the determinants defined using
functions $\veps$ from~$S_0$ to the set of edges of the graph such that
$\veps(s) \in \{s_0,s_1\}$, together with the complementary function
$\tilde\veps$ such that $\{\veps(s),\tilde\veps(s)\} =
\{s_0,s_1\}$
\[
	U^{S_0} = \frac1 {2^{\#(S_0)}} \sum_{\veps} U^{\veps(S_0), \tilde\veps(S_0)}.
\]
\end{Definition}
Since the symmetry of the matrix means changing $\veps$
to $\tilde\veps$ does not change the term, $U^{S_0}$ has at most
$2^{\#(S_0)-1}$ different terms. With these objects, we can easily state our main theorem
\begin{Theorem}\label{fund}
The integrand of a massless vacuum diagram in space-time dimension~$d$ with numerator expressed
as a product of $n$ scalar products can be written as a sum over the set
$\mathcal P_k$ of partitions of~$S$ in $k$~parts, with $k$ taking all values for which this set is
non empty
\begin{gather*}
	I = \sum_{k=1}^{\#(S)} \frac{(-1)^{\#(S)-k}\Gamma\bigl(\frac d2 + k\bigr)}{U^{\frac d2 +k} }
			\sum_{\{S_1,\ldots,S_k\} \in \mathcal P_k} \prod_{j=1}^k U^{S_j}.
\end{gather*}
We can observe that this rational function has homogeneity $-L D/2 - n$.
The value of the residue of the diagram is then the projective integral obtained by multiplying $I$ with the
product~${\prod_i \frac 1{\Gamma(\beta_i)} x_i^{\beta_i-1} {\rm d }x_i^{\vphantom{\beta_i}}}$,
indexed on all edges of $\tilde G$. Since it is a projective integral, any delta function fixing a sum
of~$x_i$ give the same result. The exponents $\beta_i$ associated to the edges must sum up to $L D/2 + n$ to ensure that the integral is
projective.
\end{Theorem}


% Original source lines 274--474
\section{Reducing to the completed diagram}
\subsection{A single scalar product}

In the case of a propagator graph $G$ in a scalar theory, it is known that the residue can be expressed as an integral which
only depends on the completed graph $\tilde G$ and its graph polynomial~$\tilde U$, with an exponent
for the additional edge giving scale invariance to the completed graph, which is then a vacuum graph (see, e.g.,~\cite{Sch10})
\begin{equation} \label{InS}
	\cI_{\tilde G} [\beta_i] = \frac {\Gamma(d/2)} {\prod_i \Gamma(\beta_i)}
			\int \prod_i \bigl( {\rm d}x_i^{\vphantom {\beta_i}} x_i^{\beta_i-1}\bigr)
			\tilde U^{-d/2} \delta \bigl(\textstyle{ \sum' }x_i -1\bigr).
\end{equation}
In both products, the index~$i$ goes over all the edges of $\tilde G$, while the sum $\sum'$ is over any non
empty subset.
The derivation of formula~\eqref{InS} uses the relation between the graph polynomial~$\tilde U$ of the completed
graph and the two Symanzik polynomials of the propagator graph: the first Symanzik polynomial is what have been called the graph
polynomial up to now and will be noted~$U$, while the second one, in this single scale setting, is $p^2 V$. The
deletion contraction relation then gives that $\tilde U = x_0 U + V$. It is traditional to give the
index~0 to the edge completing the graph.

In the case of a non trivial numerator, such an expression depending only on a completed graph seems harder to get by. Let us first study the case where the numerator is simply the scalar product $p_i\cdot p_j$ and redo the same steps as in the scalar case.
The first step is to integrate on the loop momenta. For fixed Schwinger parameters, it is a Gaussian integration with two
contributions for the scalar product, one coming from the correlation between $p_i$ and~$p_j$ which gets a factor of the dimension $d$, and the other stemming from the part proportional to the exterior momentum in both objects.
This latter part can be determined by the electric circuit analogy, where the $x_i$ parameters play the roles of
the resistance of the edges. The fraction of the intensity going through a given edge is given by the ratio of two
polynomials of degree $L$, one dependent on the edge $V^i$ and the universal $U$. By this analogy, for given
values of the parameters $x_i$, the mean of the momentum $p_i$ is given by $ p V^i/U$.

The evaluation of the diagram is then given, according to~\cite{Na71}, by
\[%\label{naive}
	\cI_G(p) = \frac {1}{ \prod_k \Gamma(\beta_k)} \int \prod_k \bigl( {\rm d}x_k^{\vphantom {\beta_k}} x_k^{\beta_k-1}\bigr)
 		\left(\frac d2 \frac{U^{i,j}}{U^{d/2+1}} + \frac{V^i V^j p^2}{U^{d/2+2}} \right) \exp\left(-p^2 \frac V U\right).
\]
The two terms do not have the same homogeneity degrees, since the $V^k$ have the same degree as $U$, while $U^{i,j}$ has one degree less.
Let us define $\beta_0$ such that $\beta_0 + \sum_i \beta_i = (L+1)d/2 +1$, so that the first term is homogeneous of degree $d/2 - \beta_0$
and the second is homogeneous of degree~${d/2 - \beta_0 +1}$ and integrate on the global scale of the Schwinger parameters. The $p^2$ already present in the second term makes the two terms proportional to \smash{$\bigl(p^2\bigr)^{\beta_0-d/2}$} so that the integrand without the $\Gamma$ prefactors and powers of the $x_j$ reads
\[%\label{naive_d}
 		\frac {d}2 \frac {\Gamma(d/2-\beta_0) U^{i,j}}{U^{\beta_0+1}V^{d/2-\beta_0}} + \frac{ \Gamma(d/2-\beta_0+1)V^i V^j}{U^{\beta_0+1}V^{d/2-\beta_0+1}} .
\]
We write $d/2$ in the first term as $\beta_0 + (d/2 - \beta_0)$ and use the functional identity for the $\Gamma$--function to obtain $\beta_0 \Gamma(d/2 - \beta_0) + \Gamma(d/2 - \beta_0 +1)$. The term with the $\beta_0$ factor can be written as an integral on $x_0$ as in the scalar case, giving
\begin{equation} \label{beta+1}
	\int_{x_0 =0}^\infty {\rm d}x_0 ^{\vphantom{\beta_0}} x_0 ^{\beta_0} \frac { \beta_0 \Gamma(d/2 + 1) }{\Gamma(\beta_0+1) } \frac { U^{i,j} } { \tilde U ^{d/2 +1} }.
\end{equation}
The $\beta_0$ factor of the numerator can be used to change the argument of the $\Gamma$--function in the denominator to $\beta_0$.

The other term has a denominator with a global degree in $U$ and $V$ of $d/2 +2$ and a simple application of the same transformation would give a higher power of $\tilde U$ in the denominator. We obtain indeed
\begin{equation}\label{intermediaire}
	\Gamma(d/2 - \beta_0 +1) \frac { V U^{i,j} + V^i V^j } { U^{\beta_0+1} V^{d/2 - \beta_0 +1}}
\end{equation}
by putting everything on the same denominator. However, the Dodgson identities allow to simplify the numerator. In the completed graph, we have
\begin{equation} \label{dodV}
	\tilde U \tilde U^{0i,0j} = \tilde U^{0,0} \tilde U^{i,j} - \tilde U^{0,j} \tilde U ^{i,0}.
\end{equation}
The part proportional to $x_0$ is trivially true but the part independent of~$x_0$ is quite interesting. It is obtained by adding the subscript 0
to the polynomials which do not have it in their upper indices. The definition of $\tilde U$ gives us that $\tilde U^{0,0} = U$
and~$\tilde U_0 = V$. The polynomial $V^k$ is defined by considering the trees in~$G$ including the edge $k$, with a sign dependent on the
orientation of~$k$ with respect to the path going from the entrance to the exit of the diagram~$G$. In the case of $\tilde U^{0,k}$, we must
consider subgraphs with a unique loop including both $0$ and~$k$, and a~sign is added if the two do not have the same orientation with respect to
this loop. The two sets of subgraphs are clearly in correspondence, with matching determination of the signs, so that we have that $\tilde U^{0,k}
= \tilde U^{k,0} = V^k$. The part for $x_0=0$ of equation~\eqref{dodV} can therefore be written~$V U^{i,j} = U \tilde U^{i,j}_0 - V^i V^j$.
The numerator in equation~\eqref{intermediaire} turns out to be proportional to~$U$, reducing the total power of $U$ and
$V$ in the denominator. The trick of writing this term as an integral over $x_0$ will give a~contribution similar to the one in equation~\eqref{beta+1}, so that both combine to give
\[%\label {fin}
	\int_{x_0 =0}^\infty {\rm d}x_0 ^{\vphantom{\beta_0}} x_0 ^{\beta_0-1} \frac { \Gamma(d/2 + 1) }{\Gamma(\beta_0)} \frac { x_0 U^{i,j} + \tilde U^{i,j}_0 } { \tilde U^{d/2+1} },
\]
where it is easy to recognize the full $\tilde U^{i,j}$ in the numerator.


\subsection{The general case}	\label{reducgene}

In the general case, such a step by step reduction of the terms coming from the parts proportional to the
exterior momentum would become rapidly very complex, without a clear generalization for an arbitrary
numerator. Indeed, when a numerator is a product, one must generate independently all possible powers of
$x_0$, splitting some terms in numerous parts. It is therefore impossible to separate a phase where one
reduces to the vacuum graph and one where the correlations are expressed in terms of the polynomials $U^S$.

However, it is possible to circumvent such a painful computation. We just start from the remark that
in a massless theory, the dependence on the exterior momentum is fixed by homogeneity reason to be a
power of $p^2$ which will be written $\beta_0 - d/2$. If we integrate $\exp\bigl(-p^2\bigr) \cI_G(p)$ on the
whole momentum space, one gets \smash{$ \Gamma(\beta_0)\pi^{d/2}/\Gamma(d/2) $} times the residue of the
graph, with the $d$ dependent factor simply half the volume of the $(d-1)$-sphere. This is linked to the use of~$p^2$ as the radial variable.
On the other hand, we could do the integration over this exterior~$p$ before all parametric
integrations, alongside the integrations on the momenta on the edges of the graph. These integration steps are just, for fixed Schwinger
parameters, the integration on the
completed graph $\tilde G$.
The parameter $x_0$ associated to the additional edge is simply fixed to~1. This constraint on $x_0$ can be
seen as a special case of the reduction of a projective integral if we complete the integrand
by a factor that makes it scale invariant. The relation between the original diagram $G$ and the completed one~$\tilde
G$ with one more loop shows that the required factor is $x_0^{\beta_0-1}$. The $\Gamma$ factors in the relation
between the residue of~$G$ and the integral on~$\tilde G$ ensure that all edges of the completed diagram are associated
with integrals with the same factor $1/\Gamma(\beta_i)$. The case of the propagator
diagrams is therefore reduced to the one of vacuum diagrams.

A similar approach is also possible in the case of non-scalar propagators. Lorentz invariance and
the presence of a single available vector give a limited number of tensor structures. In the case of a vector propagator, the only two possible
Lorentz invariant tensor structure are~$g_{\mu\nu}$ and~$p_\mu p_\nu$. The completion to a vacuum graph using these same two structures
would give enough information from these two vacuum graphs to fully determine the propagator.

The case of massive theories is less clear cut. The scale invariance which allowed this simple derivation is no longer available, but it would
nevertheless be useful to reduce the maximal power in the denominator. The step by step derivation of the previous subsection would be however
difficult to generalise: I would rather bet on careful expansion of an analog of our final formula to prove that they allow to recover Nakanishi's
results from a simpler expression. This would however require its own exposition space.

\section{Proof of the general result}
The proof of Theorem~\ref{fund} is based on the usual expression of the integral of polynomials under a~Gaussian measure as a sum of products of elementary
correlators. In a first step, the different terms are grouped in a slightly redundant manner involving permutations of the $n$ scalar product in
the set~$S$. The second step show that the same expression is obtained from the formula in Theorem~\ref{fund} through the use of Dodgson identities and
simple combinatorial identities.

\subsection{The terms to find}

From the preceding section, only correlations in the numerator have to be considered, since the studied cases
can be reduced to vacuum diagrams. The parametric method consists in replacing the propagators by integrals on parameters of Gaussian functions of
the momenta. Mix in additional variables to write the conservation of momentum at each vertex as just an additional quadratic term in the
exponential, and for fixed Schwinger parameters, a Gaussian integral has to be evaluated. In the absence of numerators and for vacuum diagrams, one
just obtains a~power of the determinant of the quadratic form in the Gaussian, which is just the first Symanzik polynomial. The numerators envisaged
in Theorem~\ref{fund} are just polynomial of degree~$2n$ in the momenta, part of the variables in the Gaussian integral. It is well known that the integral
of such polynomials reduce to a sum of products of $n$ elementary correlations of the Gaussian measure and that there are
$ (2n-1)!! $ terms in the sum, corresponding to the possible groupings by pairs of the $2n$ factors in the polynomial.\footnote{Our main interest
is with the case where all these vectors are associated to different edges of the graph, but having many to one relations between the vectors and the
edges of the graph does not change our proof. However, writing explicitly this function would make these already cumbersome formulas illegible.}
If we name $\cal T$ the
set of possible pairings~$T$ of the~$2n$ vectors in the scalar products, the numerator multiplies the gaussian integral by the following factor
\begin{equation} \label{numerat}
	\sum_{T \in \cal T} \prod_{\{i,j\} \in S} g_{\mu_i,\mu_j} \prod_{\{k,l\} \in T} g^{\mu_k,\mu_l} \frac{U^{k,l}}{2U}.
\end{equation}
In each term, the $n$ factors of the metric in the product over $S$ and the $n$ factors of its inverse in the product over $T$ combine to
only produce a power of the dimension~$d$. In order to prove Theorem~\ref{fund}, the different partitions $T$ in $\cal T$ must be collected in groups
giving the same power of~$d$. The power of $d$ is the number of independent chains of contracted indices between alternating metrics (coming
from~$S$) and inverse metrics (coming from~$T$).

The terms in Theorem~\ref{fund} are characterized by a permutation $\sigma$ of $S$ and the couple of functions~$(\varepsilon, \tilde\varepsilon)$,
as will be precised in the next subsection. Going from such data to a partition in pairs $T$ can be done by taking the pairs in~$T$ to be of the
form $\{\varepsilon(i), \tilde\varepsilon(\sigma(i)) \} $. It is clear that this produces a partition of the $2n$ vectors in $n$ pairs. In the
reverse direction, the permutation can be built cycle by cycle. Starting from a pair in~$T$ which is not already covered, we must choose the
element which will be $\varepsilon(i)$. The second element of the pair then determines what is~$\sigma(j)$ as well as the value
of~$\tilde\varepsilon$ at the point~$\sigma(j)$ if it is different from~$j$. Starting from~$\varepsilon(\sigma(j))$, the same reasoning
allows to determine successively the $\sigma^k(j)$ until one comes back to~$j$.

For a given term, the permutation $\sigma$ as well as the choice of functions $(\varepsilon,\tilde
\varepsilon)$ are not unique: in the case where $\sigma$ is the identity permutation, the product will
not depend on $(\varepsilon,\tilde \varepsilon)$ and in any case, exchanging the two functions
$(\varepsilon,\tilde \varepsilon)$ and changing $\sigma$ to its inverse leave the term unchanged. More
generally, the number of different choices of the function $\varepsilon$ compatible with a given term is
equal to $2^c$, where $c$ is the number of cycles of the permutation $\sigma$. This can be seen from the derivation of $\sigma$ from~$T$,
where we have a binary choice to make for each cycle.

When taking the average over the $2^n$ possible values of the pair $(\varepsilon,\tilde \varepsilon)$, this gives the product of the 2 in each of
the denominators for the correlations in equation~\eqref{numerat}. However, in this scheme each terms is produced multiple times and must be divided
by~$2^c$ to compensate if the permutation $\sigma$ has $c$ cycles. This number of cycles
also appears in the dependence on the dimension: inverse metrics from the correlations and metrics from the scalar products combine to give the trace of
the identity on the momentum space for each cycle of~$\sigma$, giving a factor~$d^c$.
Combining these two factors, the terms with a permutation of $c$ cycles must be multiplied by~$\bigl(\frac d 2 \bigr)^c$.
The signature of a permutation has also a simple expression in terms of the number of cycles, it is simply
the parity of $n-c$, with $n$ the number of objects on which the permutation acts.

\subsection[The terms in U\^{}S\_0]{The terms in $\boldsymbol{U^{S_0}}$}

We now start from the expression appearing in Theorem~\ref{fund}. Since all terms are subject to the
same averaging on all the possible values of the $\varepsilon$ function, we can just forget about the
different maps $\varepsilon$ and consider a fixed one. If we put every term on the same denominator
$U^{d/2+ \#(S)}$, each $U^{\varepsilon(S_i),\tilde\varepsilon(S_i)}$ can receive a factor $U^{\#(S_i) - 1}$
and be converted by the Dodgson identities
into the determinant of a matrix with correlators $U^{i,j}$ as elements.

If we now expand each of these determinants in terms of permutations, we obtain terms which are entirely
similar to the ones obtained in the preceding section. A given permutation of the whole set~$S$ can be
obtained in a number of different ways according to its cycle structure. Indeed, a given cycle structure can
only be obtained from the product of determinants associated to a partition of~$S$ in subsets which are
union of (the sets subjacent to) cycles. The identity is for example compatible with any partition of~$S$,
while a permutation with only one cycle is only compatible with the trivial partition with only one subset.

For a given permutation, the signs coming from the various determinants always combine to the signature of the complete
permutation, since the signature is multiplicative. It is given by the
parity of the size of~$S$ minus the number of cycles. Combined with the sign in our formulas, we obtain
that the unique term with the finest partition compatible with the cycle structure will have a positive sign
and the factor $\Gamma(d/2 + c)$, with $c$ the number of cycles in the permutation. We therefore will
have contributions from partitions of~$S$ with~$c$ down to $1$~parts, which comes with alternating
signs. All terms with the same number of parts in the partition are identical. Their numbers are
given by the Stirling numbers of the second kind $\left\{\begin{smallmatrix} c \\ p \end{smallmatrix} \right\}$, which
count the number of ways the $c$ cycles of the permutation can be collected in $p$
parts.\footnote{The notations used are the one proposed by Donald Knuth
in~\cite{Knuth_1992}. They have the advantage of circumventing the notation clash between combinatorists and analysts
which use the notation $(x)_n$ for respectively falling and raising factorials.}
We therefore obtain the following coefficients for the terms with $c$ cycles
\begin{equation}
\label{sumOfPartitions}
	\sum_{p=1}^c (-1)^{c-p}
	\left\{ \begin{matrix} c \\ p \end{matrix} \right\}
	\Gamma(d/2 + p).
\end{equation}
Now $\Gamma(d/2 + p)$ can be converted, using the functional equation for $\Gamma$, in the product of
$\Gamma(d/2)$ and the raising factorial or Pochhammer symbol $ (d/2)^{\overline p}$. It is a well-known property of the
Stirling numbers of the second kind that they allow to convert from falling factorials to plain powers, as it was
the reason for their introduction by Stirling in~\cite{Stirling_1730}, or by a
simple change of signs, from raising factorials to plain powers: all we need is the product rule
\[x \times
x^{\overline p} = x^{\overline{p+1}} - p x^{\overline p}
\]
 and the corresponding one for $x \times
x^{\underline p}$ and the triangle property of the Stirling number to obtain that the sum in
equation~\eqref{sumOfPartitions} will give $(d/2)^c \Gamma(d/2)$. This completes our proof of Theorem~\ref{fund}.

\appendix\setcounter{section}{0}
% Original source lines 1287--1327
\section{Dodgson identities}\label{polyn}

In order to obtain exact results, we need the Dogdson identities in a form which allows to track signs.
They are indeed essential to obtain exact results. In other contexts, the signs can be irrelevant
since the Dodgson polynomials will appear squared or as part of expressions the sign of which is not important,
but it is not the case here. Nevertheless, we want to keep things simple and with
signs which depend on a minimum number of choices in the presentation of the diagram. In particular, the usual
presentation as the determinant of a submatrix with rows and columns removed keeps a dependence on the order
of the rows. The convention we proposed in Section~\ref{presentation} for $U^{I ,J}$ is the determinant of the matrix
$M_G$ with the columns indexed by the set $I$ replaced with ones with only a 1 in the position indexed by the corresponding
element of~$J$. Independence on the order of the edges or the vertices is clear in this case and expansion according to
the columns indexed by~$I$ trivially gives back the usual definition up to a sign, while multilinearity can be used to
show that the same result can be obtained by changing instead the rows indexed by~$J$. Finally, in the case where $I$
and~$J$ have only one element, $U^{i,j}$ only depends on the orientations of $i$ and~$j$.

The sign introduced can be deduced as follows: the elements of $I$ have first to be brought to the beginning of the matrix,
requiring $i_1 - 1$ transpositions to send the first element of $I$ to~1,~${i_2 -2}$ for the following up to $i_k -k $
for the last one if we suppose that the elements of~$I$ are numbered according to the ordering of edges and not the order in~$I$.
We have similar number of transpositions for bringing the elements of $J$ to the front, and the total sign only depends on the parity
of the sum $\sum_{i\in I} i+ \sum_{j\in J} j$. The upper left block of the resulting matrix is a~permutation matrix, which gives a sign
corresponding to the product of the signs associated to the permutations bringing $I$ and~$J$ in the natural order. If one recalls that
the convention in this work is to put a minus sign before one of the incidence block rather than to use the symmetric form used by Schnetz,
one recognizes that there are all the elements of the sign proposed in \cite[equation~(8)]{SC2021}.

Our aim is to express some determinants involving correlators~$U^{i,j}$. Examining the proof of Dodgson
identities, such
as the one found in~\cite{Br2009}, shows that our definition of~$U^{I ,J}$ naturally appears in them. Indeed, $U^{i,j}$
is the $j,i$ component of the adjugate matrix of $M_G$, the matrix which, multiplied by $M_G$ gives $\det(M_G)$
times the identity. We want to express the determinant of the matrix $M$ with elements $U^{i,j}$ such that $i$ is in $I$ and $j$ is in
$J$. We form a matrix of the dimension of~$M_G$ by replacing in the identity matrix the columns with index in~$J$ by the
ones indexed by~$I$ of the adjugate matrix. An expansion along the columns coming from the identity matrix shows
that this new matrix has the same determinant as~$M$.

Multiplying by the left by~$M_G$, one obtains a matrix which has the columns of~$M_G$ except for the columns indexed
by~$J$ which have $\det(M_G)$ in the position indexed by~$I$. We recognize the matrix used to define $U^{J,I}$,
apart from the factors $\det(M_G)$, so that we obtain the required Dodgson identities without any sign ambiguity by
taking the determinants
\[
	\det(M_G) \det\bigl( (U^{i,j})_{i\in I,j\in J} \bigr) = \det(M_G)^{\# J} U^{I ,J}.
\]

\begin{thebibliography}{99}
\bibitem[3]{Br2009}
Brown F., On the periods of some {F}eynman integrals, \href{https://arxiv.org/abs/0910.0114}{arXiv:0910.0114}.

\bibitem[6]{BrYe09}
Brown F., Yeats K., Spanning forest polynomials and the transcendental weight
 of {F}eynman graphs, \href{https://doi.org/10.1007/s00220-010-1145-1}{\textit{Comm. Math. Phys.}} \textbf{301} (2011),
 357--382, \href{https://arxiv.org/abs/0910.5429}{arXiv:0910.5429}.

\bibitem[12]{Knuth_1992}
Knuth D.E., Two notes on notation, \href{https://doi.org/10.2307/2325085}{\textit{Amer. Math. Monthly}} \textbf{99}
 (1992), 403--422, \href{https://arxiv.org/abs/math.HO/9205211}{arXiv:math.HO/9205211}.

\bibitem[16]{Na71}
Nakanishi N., Graph theory and {F}eynman integrals, Vol.~11, Gordon and Breach,
 New York, 1971.

\bibitem[18]{Sch10}
Schnetz O., Quantum periods: {A}~census of {$\phi^4$}-transcendentals,
 \href{https://doi.org/10.4310/CNTP.2010.v4.n1.a1}{\textit{Commun. Number Theory Phys.}} \textbf{4} (2010), 1--47,
 \href{https://arxiv.org/abs/0801.2856}{arXiv:0801.2856}.

\bibitem[19]{SC2021}
Schnetz O., Geometries in perturbative quantum field theory, \href{https://doi.org/10.4310/CNTP.2021.v15.n4.a2}{\textit{Commun.
 Number Theory Phys.}} \textbf{15} (2021), 743--791, \href{https://arxiv.org/abs/1905.08083}{arXiv:1905.08083}.

\bibitem[20]{Stirling_1730}
Stirling J., Methodus differentialis: sive tractatus de summatione et
 interpolatione serierum infinitarum, Londini, Typis Gul. Bowyer, 1730,
 available at \url{http://archive.org/details/bub_gb_71ZHAAAAYAAJ}.

\end{thebibliography}
\end{document}
